\documentclass{ximera}

\begin{document}

    \title{Lecture 21}
    
    \begin{abstract}  
    Outline: \\
    - Rolle's Theorem \\
    - Mean Value Theorem
\end{abstract}

\maketitle

\begin{theorem}{(Rolle's Theorem)}
    Let $f:[a,b] \to \mathbb{R}$ be continuous on $[a,b]$ and differentiable on $(a,b)$.  If $f(a) = f(b)$, then there exists a point $c \in (a,b)$ where $f'(c) = 0$.
\end{theorem}
\begin{proof}
    Since $f$ is continuous on a compact set, it attains a max and a min due to EVT.  If they both occur at the endpoints, then $f$ is a constant function and consequently $f'(x) = 0$ for all $x \in [a,b]$.  Else, if the max or min occurs on the interior, then it follow from the Interior Extremum Theorem that there exists a $c \in (a,b)$ where $f'(c) = 0.$
\end{proof}

\begin{theorem}{(Mean Value Theorem}
    If $f:[a,b]\to \mathbb{R}$ is continuous on $[a,b]$ and differentiable on $(a,b)$, then there exists a point $c \in (a,b)$ where 
    \begin{align*}
        f'(c) = \frac{f(b) - f(a)}{b-a}.
    \end{align*}
\end{theorem}
\begin{proof}
    Notice this is a more general version of Rolle's Theorem when $f(b) = f(a)$.  To use Rolle's Theorem, we consider the line
    \begin{align*}
        f(a) + \left(\frac{f(b) - f(a)}{b-a} \right)(x-a).
    \end{align*}
    Now consider
    \begin{align*}
        d(x) = f(x) - \left(f(a) + \left(\frac{f(b) - f(a)}{b-a} \right)(x-a)\right).
    \end{align*}
    Then, $d(a) = f(a) - f(a) = 0$, and $d(b) = f(b) -\left(f(a) + \left(\frac{f(b) - f(a)}{b-a} \right)(b-a)\right) = 0$.
    Consequently, we can apply Rolle's Theorem to the function $d(x)$:
    there exists a $c \in (a,b)$ such that $d'(c) = 0$, or 
    \begin{align*}
        f'(c)  - \frac{f(b) - f(a)}{b-a} = 0.
    \end{align*}
\end{proof}

\end{document}
